\section{Dataset and Implementation}

This section details the dataset used for training and evaluation, along with specific implementation aspects of the ExplainableCropAI system.

\subsection{Dataset Description}

We utilized the Paddy Doctor dataset \cite{das2020paddy}, a comprehensive collection of paddy disease images specifically designed for automated disease classification research. The dataset contains 10,407 high-resolution RGB images (3264 $\times$ 2448 pixels) captured in real field conditions across different growing regions in South and Southeast Asia.

\begin{table}[h]
\centering
\caption{Class Distribution in the Paddy Doctor Dataset}
\label{tab:dataset-distribution}
\begin{tabular}{lcc}
\hline
\textbf{Class} & \textbf{Images} & \textbf{Percentage} \\
\hline
Bacterial Leaf Blight & 1,584 & 15.2\% \\
Brown Spot & 1,676 & 16.1\% \\
Leaf Blast & 1,738 & 16.7\% \\
Bacterial Leaf Streak & 1,237 & 11.9\% \\
Sheath Blight & 1,056 & 10.1\% \\
Healthy & 1,488 & 14.3\% \\
Other Disease 1 & 412 & 4.0\% \\
Other Disease 2 & 373 & 3.6\% \\
Other Disease 3 & 427 & 4.1\% \\
Other Disease 4 & 416 & 4.0\% \\
\hline
\textbf{Total} & \textbf{10,407} & \textbf{100\%} \\
\hline
\end{tabular}
\end{table}

The dataset presents several challenges characteristic of real-world agricultural scenarios:

\begin{itemize}
    \item \textbf{Variable lighting conditions}: Images captured under different natural lighting conditions throughout the day
    \item \textbf{Background complexity}: Presence of soil, water, and other plants in many images
    \item \textbf{Disease co-occurrence}: Some plants exhibiting multiple diseases simultaneously
    \item \textbf{Disease progression}: Varying severity levels from early to advanced stages
    \item \textbf{Image quality variation}: Some images affected by blur, glare, or shadow
\end{itemize}

\subsubsection{Data Preprocessing and Augmentation}

We preprocessed the images using the following pipeline:

\begin{enumerate}
    \item Resizing to 256 $\times$ 256 pixels to standardize input dimensions
    \item Normalization using ImageNet mean and standard deviation values (mean=[0.485, 0.456, 0.406], std=[0.229, 0.224, 0.225])
    \item Center cropping to 224 $\times$ 224 pixels during inference
\end{enumerate}

For training, we applied data augmentation techniques to improve model robustness:

\begin{itemize}
    \item Random horizontal and vertical flips (probability 0.5)
    \item Random rotation ($\pm$15 degrees)
    \item Random brightness and contrast adjustments ($\pm$10\%)
    \item Random cropping (224 $\times$ 224 pixels)
    \item Mixup augmentation \cite{zhang2018mixup} with $\alpha=0.2$
\end{itemize}

\subsubsection{Dataset Splits}

We divided the dataset into training (70\%), validation (15\%), and testing (15\%) sets using stratified sampling to maintain class distribution. To ensure robust evaluation, we created three evaluation subsets:

\begin{enumerate}
    \item \textbf{Standard test set}: Random stratified sample from the main dataset
    \item \textbf{Challenging test set}: 500 images with characteristics that typically challenge classifiers (low light, unusual angles, early disease stages)
    \item \textbf{Explanation evaluation set}: 500 images with pixel-level annotations of disease-affected regions for evaluating explanation quality
\end{enumerate}

\subsection{Implementation Details}

\subsubsection{Model Architecture}

Our base disease detection model was implemented using PyTorch 1.9.0. We used ResNet50 pretrained on ImageNet as the backbone, with the following modifications:

\begin{lstlisting}[language=Python, caption=Modified ResNet50 for explainability, label=lst:resnet-model]
class PaddyDiseaseClassifier(nn.Module):
    def __init__(self, num_classes=10, pretrained=True):
        super(PaddyDiseaseClassifier, self).__init__()
        
        # Load pretrained ResNet50 without the classification head
        self.backbone = models.resnet50(pretrained=pretrained)
        self.backbone = nn.Sequential(*list(self.backbone.children())[:-2])
        
        # Global Average Pooling
        self.gap = nn.AdaptiveAvgPool2d(1)
        
        # Dropout for uncertainty estimation
        self.dropout = nn.Dropout(p=0.5)
        
        # Classification head
        self.fc = nn.Linear(2048, num_classes)
        
        # For storing activations and gradients
        self.gradients = None
        self.activations = None
    
    def activations_hook(self, grad):
        self.gradients = grad
    
    def forward(self, x):
        # Extract features
        x = self.backbone(x)
        
        # Save activations for Grad-CAM
        self.activations = x
        
        # Register hook for gradient capture
        if self.training or x.requires_grad:
            h = x.register_hook(self.activations_hook)
        
        # Global average pooling
        x = self.gap(x)
        x = x.view(x.size(0), -1)
        
        # Apply dropout (enabled during inference for MC Dropout)
        x = self.dropout(x)
        
        # Classification
        logits = self.fc(x)
        
        return logits, {'activations': self.activations}
\end{lstlisting}

\subsubsection{Training Procedure}

The model was trained with the following hyperparameters:

\begin{itemize}
    \item Optimizer: Adam with learning rate 0.0001
    \item Learning rate scheduler: ReduceLROnPlateau (patience=5, factor=0.1)
    \item Loss function: Cross-entropy loss
    \item Batch size: 32
    \item Epochs: 100 (with early stopping, patience=10)
    \item Weight decay: 0.0001
\end{itemize}

Training was performed on an NVIDIA Tesla V100 GPU with 16GB memory. The final model was selected based on validation accuracy, achieving 97.5\% on the validation set and 97.2\% on the standard test set.

\subsubsection{Visual Explanation Implementation}

The Grad-CAM implementation follows Algorithm \ref{alg:gradcam}, which generates class-discriminative localization maps using the gradients of the target class flowing into the final convolutional layer.

\begin{algorithm}
\caption{Gradient-weighted Class Activation Mapping (Grad-CAM)}
\label{alg:gradcam}
\begin{algorithmic}[1]
\REQUIRE Image $I$, CNN model $f$, Target class $c$, Final convolutional layer $l$
\ENSURE Class activation map $L_{Grad-CAM}^c$

\STATE $y^c \gets f(I)_c$ \COMMENT{Forward pass to get score for class $c$}
\STATE $\frac{\partial y^c}{\partial A^l} \gets$ Compute gradients of $y^c$ w.r.t. feature maps $A^l$
\STATE $\alpha_k^c \gets \frac{1}{Z}\sum_i\sum_j \frac{\partial y^c}{\partial A_{ij}^k}$ \COMMENT{Global average pooling of gradients}
\STATE $L_{Grad-CAM}^c \gets ReLU\left(\sum_k \alpha_k^c A^k\right)$ \COMMENT{Weighted combination of feature maps}
\STATE Upsample $L_{Grad-CAM}^c$ to input resolution
\RETURN $L_{Grad-CAM}^c$
\end{algorithmic}
\end{algorithm}

\subsubsection{Uncertainty Quantification Implementation}

For uncertainty quantification, we implemented Monte Carlo Dropout by enabling dropout during inference and performing 30 forward passes for each input image. The uncertainty metrics were calculated as follows:

\begin{lstlisting}[language=Python, caption=Uncertainty estimation with Monte Carlo Dropout, label=lst:mc-dropout]
def estimate_uncertainty(model, image, num_samples=30):
    predictions = []
    
    # Enable dropout during inference
    model.train()
    
    # Multiple forward passes
    for _ in range(num_samples):
        with torch.no_grad():
            logits, _ = model(image)
            probs = F.softmax(logits, dim=1)
            predictions.append(probs.cpu().numpy())
    
    # Stack predictions
    predictions = np.stack(predictions)  # (num_samples, batch_size, num_classes)
    
    # Mean prediction
    mean_probs = np.mean(predictions, axis=0)
    
    # Predictive entropy (total uncertainty)
    entropy = -np.sum(mean_probs * np.log(mean_probs + 1e-10), axis=1)
    
    # Expected entropy (aleatoric uncertainty)
    expected_entropy = -np.mean(
        np.sum(predictions * np.log(predictions + 1e-10), axis=2),
        axis=0
    )
    
    # Mutual information (epistemic uncertainty)
    mutual_information = entropy - expected_entropy
    
    # Class agreement
    pred_classes = np.argmax(predictions, axis=2)
    most_common = scipy.stats.mode(pred_classes, axis=0)[0]
    agreement = np.mean(pred_classes == most_common, axis=0)
    
    return {
        'entropy': entropy,
        'mutual_information': mutual_information,
        'expected_entropy': expected_entropy,
        'agreement': agreement,
        'mean_probs': mean_probs
    }
\end{lstlisting}

\subsubsection{Natural Language Explanation Templates}

The natural language explanation generator uses a template-based approach with domain-specific knowledge. Templates were developed in collaboration with agricultural experts and include:

\begin{itemize}
    \item Disease detection templates (e.g., "This image shows a paddy plant with \{disease\}")
    \item Visual evidence templates (e.g., "The key visual indicators include \{features\}")
    \item Confidence templates at three levels (high, medium, low)
    \item Treatment recommendation templates
    \item Prevention advice templates
    \item Comparison templates for similar diseases
\end{itemize}

The templates are populated with disease-specific information from the knowledge base and combined based on the prediction confidence and detected visual features.

\subsubsection{Recommendation Engine Implementation}

The recommendation engine was implemented as a rule-based system using a structured knowledge base of disease information and treatment protocols. The knowledge base includes:

\begin{itemize}
    \item Disease symptoms and progression patterns
    \item Treatment options at different severity levels
    \item Prevention strategies
    \item Economic thresholds for intervention
    \item Environmental risk factors
\end{itemize}

Recommendations are generated by combining the disease prediction, uncertainty metrics, and additional contextual information such as crop stage and environmental conditions.

\subsection{System Integration}

The individual components were integrated into a unified framework using a modular architecture. The system flow follows these steps:

\begin{enumerate}
    \item The image is processed by the base model to generate a disease prediction
    \item Grad-CAM is applied to generate visual explanations
    \item Monte Carlo Dropout is performed to estimate uncertainty
    \item Visual features are extracted from the explanation heatmap
    \item Natural language explanations are generated using the templates
    \item Recommendations are produced based on the prediction, uncertainty, and context
    \item The complete output (prediction, visual explanation, textual explanation, and recommendations) is presented to the user
\end{enumerate}

\subsection{User Interface}

The user interface was designed to present the multi-level explanations in an intuitive manner. Figure \ref{fig:user-interface} shows the interface design with key components:

\begin{enumerate}
    \item Original image with disease prediction and confidence
    \item Grad-CAM visualization highlighting important regions
    \item Natural language explanation of the diagnosis
    \item Uncertainty visualization showing confidence distribution
    \item Treatment and prevention recommendations
    \item Option to provide additional context for more tailored recommendations
\end{enumerate}

\begin{figure}[ht]
    \centering
    % [Insert UI screenshot here]
    \caption{ExplainableCropAI user interface showing integrated prediction, explanation, and recommendation components.}
    \label{fig:user-interface}
\end{figure}

The interface was implemented as a responsive web application and a mobile app to support various deployment scenarios, from laboratory settings to field use on mobile devices.
